\documentclass[UTF8,fontset=none,11pt,a4paper]{ctexart}
\usepackage{ipara-handbook}
\handbooksetup{DS-07 图的表示与遍历}{408 · 数据结构 DS-07}{Graph · Representation · Coverage}
\begin{document}
\handbooktitle{图的表示与遍历}{关系怎样被表示，并从起点覆盖而不重复失控}{408 · 数据结构 Topic DS-07}

\begin{corebox}{母问题：图的关系如何变成可验证的访问过程？}
图遍历的生成链是
\[
\text{Vertices/Edges}\to\text{Representation}\to\text{Frontier}\to\text{Visited State}\to\text{Expansion}\to\text{Coverage}.
\]
邻接矩阵用空间换常数时间的边查询，邻接表按实际边数保存关系；DFS 与 BFS 都要维护 visited，区别在 frontier 的 LIFO/FIFO 取出规则。遍历只负责覆盖和发现关系，不自动解决最短路、MST 或拓扑目标。
\end{corebox}
\begin{methodbox}{本册定位与 Owner 边界}
\begin{itemize}
\item \kw{Owns：}图、邻接矩阵/邻接表/十字链表/邻接多重表、方向与重边语义、DFS/BFS、visited、连通分量与遍历森林。
\item \kw{Uses：}DS02 的栈/队列端点契约和 DS-B01 frontier 接口；DS08 使用本册表示。
\item \kw{Stops：}最短路、最小生成树、拓扑排序等全局目标归 DS08；网络路由只 Use。
\end{itemize}
\end{methodbox}
\tableofcontents\newpage

\section{表示先于算法}
邻接矩阵 $A[i][j]$ 直接记录边存在性或权值，适合稠密图和频繁边查询，空间 $\Theta(V^2)$。邻接表为每个顶点保存出边列表，空间 $\Theta(V+E)$，遍历一个顶点的邻居只扫描实际出边。无向边需要在两端各登记一次；有向边只登记出发端。平行边和自环是否允许，必须在接口中先固定。

\subsection{十字链表：一条有向弧同时进入出边链与入边链}
普通有向邻接表让枚举顶点 $u$ 的出边很自然，但若要枚举入边，往往还需逆邻接表。十字链表把一条弧记录为
\[
(tail,head,\;tlink,\;hlink,\;info),
\]
其中 $tlink$ 连接同弧尾的下一条出边，$hlink$ 连接同弧头的下一条入边；顶点记录同时保存首条出弧和首条入弧。一个边对象被两条链引用，但只存一份权值和附加信息。

不变量是：每条弧在其尾顶点的出边链中出现一次，也在其头顶点的入边链中出现一次；删除弧必须同时从两条链摘除。它适合同时需要入度、出度和双向邻接扫描的有向图。

\subsection{邻接多重表：无向边是共享对象，而不是两份边副本}
无向邻接表通常把边 $\{u,v\}$ 分别登记在 $u,v$ 的链中。邻接多重表为每条无向边只建一个边节点，保存两个端点 $ivex,jvex$ 以及分别连接两个端点下一条关联边的 $ilink,jlink$。从顶点 $u$ 扫描时，要根据 $u$ 是该边的哪一个端点选择对应 link。

其收益是边级插删与标记只作用于一个共享对象；代价是遍历 link 时必须判断当前端点角色。十字链表面向有向弧的入/出双链，邻接多重表面向无向边的两端关联，二者不能只凭“都有两条链”互换。

\begin{methodbox}{表示选择的题面信号}
只需判断任意两点是否相邻且图较稠密，优先邻接矩阵；按实际边枚举出边，邻接表足够；有向图同时频繁枚举入边与出边，看十字链表；无向图需要对同一边做删除、标记或边级维护，看邻接多重表。选择依据是需要沿哪个关系方向访问，而不是结构名称的新旧。
\end{methodbox}

\section{visited 是覆盖不变量}
遍历的核心不变量是：每个顶点最多被“首次发现”一次，已发现但未展开的顶点位于 frontier，展开后其出边已按表示规则检查。非连通图从一个起点只能覆盖一个可达分量；要得到遍历森林，必须在所有顶点上补一层外循环。
\begin{examplebox}{同一关系，不同 frontier}
从 $0$ 出发，若先把邻居压栈，得到 DFS 的深路径；若把邻居入队，得到 BFS 的层次距离。两者都可以覆盖同一可达集，但只有无权图 BFS 的首次发现层数具有最短边数语义。
\end{examplebox}

\section{DFS、BFS 与遍历森林}
递归 DFS 隐式使用调用栈；显式 DFS 必须约定邻接顺序，否则输出序列不唯一。BFS 使用队列，入队时就标记 visited，不能等出队才标记，否则多个前驱会重复入队。对邻接表，二者时间均为 $O(V+E)$，额外空间为 $O(V)$；矩阵扫描每个顶点所有候选邻居，时间为 $O(V^2)$。

\subsection{四个动词固定遍历状态}
把任何遍历逐行代码都翻译成四个动作：
\[
\boxed{Discover\to Mark\to Enqueue/Push\to Expand}.
\]
Discover 是第一次遇到未访问邻居；Mark 把“已发现”写入 visited；Enqueue/Push 把它加入待处理边界；Expand 是将来取出后扫描邻接表。Mark 必须紧跟 Discover，因为 frontier 中的顶点已经被别的路径“预约”，不能再次加入。

\begin{examplebox}{为什么出队时标记会重复}
设 $0$ 同时连接 $1,2$，且 $1,2$ 都连接 $3$。BFS 展开 $1$ 时发现 $3$ 并入队；若尚未标记，展开 $2$ 时又会把 $3$ 入队。发现时标记则把状态从 `Unseen' 立即改为 `Discovered'，第二条边只被检查，不再制造第二份 frontier 项。
\end{examplebox}

\begin{center}\small
\begin{tabularx}{\linewidth}{L{2.1cm}L{2.4cm}YY}
\toprule
顶点状态 & 所在位置 & 已经保证什么 & 下一动作\\
\midrule
Unseen & 不在 frontier & 尚无已知发现路径 & 被邻边 Discover\\
Discovered & frontier 中 & 已有一条发现路径，visited=true & 等待取出\\
Expanded & frontier 外 & 所有邻边已扫描 & 不再展开\\
\bottomrule
\end{tabularx}
\end{center}

DFS 与 BFS 共享这三态；区别只在 Discovered 顶点按 LIFO 还是 FIFO 取出。遍历森林再增加外层动作：找最小的 Unseen 顶点作为新根，重复同一状态机。

\section{代码与测试证据}
完整实现位于 \codepath{code/ds07_graph_traversal.hpp}，测试位于 \codepath{code/ds07_graph_traversal_test.cpp}。
\lstinputlisting[language=C++,firstline=12,lastline=92,firstnumber=12,numbers=left,caption={邻接表、BFS、DFS 与可达覆盖}]{30_408/10_数据结构/07_图的表示与遍历/code/ds07_graph_traversal.hpp}
\lstinputlisting[language=C++,firstline=8,lastline=45,firstnumber=8,numbers=left,caption={方向、非连通与空图测试}]{30_408/10_数据结构/07_图的表示与遍历/code/ds07_graph_traversal_test.cpp}
\begin{lstlisting}[language=bash]
clang++ -std=c++17 -Wall -Wextra -Werror -pedantic code/ds07_graph_traversal_test.cpp -o /tmp/ds07_test
/tmp/ds07_test
\end{lstlisting}

\section{边界与概念分离}
\begin{center}\small
\begin{tabularx}{\linewidth}{L{2.5cm}C{1.8cm}C{1.8cm}YY}\toprule
操作 & 邻接表 & 邻接矩阵 & 不变量/前提 & 边界\\\midrule
遍历 & $O(V+E)$ & $O(V^2)$ & 每顶点首次发现一次 & 空图、非连通\\
边查询 & $O(\deg u)$ & $O(1)$ & 方向语义固定 & 重边/自环\\
空间 & $O(V+E)$ & $O(V^2)$ & 只保存约定关系 & 稠密/稀疏\\\bottomrule\end{tabularx}\end{center}
\begin{center}\small
\begin{tabularx}{\linewidth}{L{2.3cm}C{.35cm}L{2.3cm}YY}\toprule
概念 A&$\neq$&概念 B&判据&错因\\\midrule
DFS&$\neq$&BFS&栈 vs 队列&误写距离结论\\
首次发现&$\neq$&出队/出栈&visited 时机&重复入 frontier\\
遍历覆盖&$\neq$&最短路/MST&目标不同&把结构事实当目标证明\\\bottomrule\end{tabularx}\end{center}

\section{做题控制协议}
先标顶点、方向、重边和表示；再确定起点与邻接访问顺序。看到“访问所有可达顶点”选 DFS/BFS，看到“最少边数”才升级到 BFS 距离，看到权值优化转交 DS08。每一步写出 visited 与 frontier 的变化，最后检查是否遗漏非连通分量。
\begin{corebox}{压缩复原}
五问：边存在哪里？当前 frontier 是栈还是队列？何时标记 visited？起点能覆盖哪些分量？题目要覆盖还是要优化？
\end{corebox}
\end{document}
